
In the last decade, \glspl{lib} have consolidated their prime position as the go to battery for all consumer electronics \cite{Zubi.2018}. With the climate protection targets and the demand for a rising amount of electric vehicles for private and public transport, the \gls{lib} is also the only technology that has the needed energy density, as well as the development level necessary, to support the transition from fossil fuels to \ce{CO2} free transportation. While there are new technologies developed with even higher theoretical energy densities, none of the candidates has reached a technology readiness level yet that could challenge the usability of the \gls{lib} \cite{Placke.2017}. One big concern with batteries in general however is the ongoing degradation through battery ageing during operation that not only increases the running of costs of all \glspl{ev} but also reduces the customer acceptance of the vehicles, which is vital for a widespread change to electric transportation \cite{Agubra.2013}. With increasing ageing, not only the capacity of a cell fades but also the available power. As of time the main efforts in the battery research are targeted towards increasing the energy density to enable higher range on full electric vehicles and to reduce the ageing on the cell to increase the overall lifetime \cite{Vetter.2005,Lin.2015}. One of the strongest ageing effects possible in \gls{lib}s is the so call lithium plating, which describes the unwanted deposition of metallic lithium on the anode surface during charging. This deposition takes place when the device is charged at very low temperatures, or when very high charging currents are applied. The capacity losses of plating can lead to the end of life of a cell within 10s of cycles, when the nominal cell life can last for hundreds or even thousands of cycles \cite{JiangFan.2006}. The plating process can also cause the growth of lithium dendrites, small lithium metal structures that can pierce the separator and cause shortcuts within the cell. These can lead to cell failure or in the worst case to a thermal runaway of the cell, a process where the cell reaction is self-energizing and causing a cell fire \cite{Waldmann.2018}.  \par
The conditions for lithium plating are rarely met for cells within consumer electronics, since they are normally charged at moderate rates within living areas. This is different for \gls{ev}s however. Electric cars are supposed to work under all normal outside conditions, charging at low temperatures is something that cannot be prevented. Today the vast majority of customers with electric cars also own a house or have access to a garage to charge their cars at higher temperatures under better \cite{Nobis.2019}. With the target of a much higher percentage of \gls{ev}s in the future however, this cannot be guaranteed for each owner of an electric car and some will have to be charged at public charging stations outside. To compensate for this effect, it is possible to equip \gls{ev}s with a heating system that heats up the battery to the optimal charging temperature prior to the charging process. The drawback of this method is for one the energy losses induced by the necessary heating process that increase the running costs of the car and the increased charging time due to the added heating time. In some cases there may not even be enough energy left in the cell to heat up properly to prevent plating during charging. The batteries are heating up due to the ohmic losses during the charging process, but this does not fully prevent the ageing caused by plating. Another aspect that is of vital importance for \gls{ev}s is the charging speed \cite{Farjah.2018}. To compete with combustion engine cars the charging speeds of EVs would have to be in the range of minutes. Achieving this would greatly increase the acceptance of the \gls{ev}s as a fully capable replacement. This would however require charging speed that greatly exceed the recommended nominal charging speeds for standard commercials \gls{lib}s. To fully charge a battery within ten minutes a charging speed of 6 times higher than average expected data sheet speeds is required. It is possible to manufacture cells that can handle higher charging speeds, but this is always at the cost of energy density, which is vital for the range of the \gls{ev}. For these very high charging speeds, plating effects are to be expected even at moderate temperatures, leading to very fast cell degradation \cite{Waldmann.2018}. \par
Grid stabilization using \gls{ev}s could be the solution for the strongly fluctuating energy output of renewable energy sources. All \gls{ev}s that are not currently in use and that are connected to the power grid could be used to stabilize the grid and store excess energy in times of overproduction, while at the same time providing energy in times of high demand \cite{Nunes.2017,Pramana.2017,Quinn.2010,Yan.2014}. If the vehicles are used in this way, low environment temperatures could be very damaging for the battery cells. In this case, the energy influx is unpredictable and uncontrollable and thus no prior heating is possible. Always keeping the batteries at a safe temperature however would produce very high energy losses. \par
The amount of available spent batteries is going to increase significantly within the next few years when the first generations of car batteries will reach the end of life. The plating effect is here also attributed to sudden cell death at very low state of health \cite{Waldmann.2018}, when the surface electrolyte interface (SEI) layer is getting very thick. Predicting this mode of cell failure could be very beneficial for second life usage of spent \gls{ev} batteries. \par
Understanding and preventing plating could not only increase the general lifetime for \gls{lib}, but also enable the usage of \gls{ev}s as a grid support for renewable energies, greatly increasing the usability of fully renewable energy sources to reduce \ce{CO2} emissions. It also can help make second life applications more reliable and thus extend the commercial value of the batteries. At the same time, it increases the safety of \gls{lib}s by preventing dendrites that can cause cell failure or dangerous thermal runaways.
A deeper understanding of the plating process can also help to increase the usability of metallic lithium in future battery technologies. The usage of metallic lithium as the anode in a lithium battery drastically increases the energy density. The re-deposition of the lithium on the anode during the charging process in the battery however is hard to control and leads to very low cycles life of those cells \cite{Lin.2017,Cheng.2017}. Even for completely new technologies, the usage of lithium is highly likely \cite{Placke.2017}. lithium is not only the lightest solid material in existence but it also has the highest negative electro potential. This enables very lightweight batteries with a high voltage and thus high energy density.\footnote{Footnote example} \gls{x}
